\answerpaper{非数学 A 类}

\answerproblem{一、}{30}
填空题：
\begin{enumerate}[label=(\arabic*)]
  \item 设函数 $f(x)=x-\ln(1+x)-\frac{1}{2}x\sin x$，$g(x)=ax^3$，且满足 $\displaystyle \lim_{x\to 0}\frac{f(x)}{g(x)}=1$，则 $a=\underline{\hspace{3.5em}}$．
  \item 设平面过原点与点 $(6,-3,2)$，且与平面 $4x-y+2z-8=0$ 垂直，则此平面方程为 $\underline{\hspace{3.5em}}$．
  \item 极限 $\displaystyle \lim_{x\to 0}\frac{\displaystyle \int_0^{\sin^2 x}\ln(1+t)\dd t}{(\sqrt[3]{1+x^3}-1)\sin x}=\underline{\hspace{3.5em}}$．
  \item 设函数 $f(u,v)$ 具有一阶连续偏导数，$f(2,0)=3$，$f_u'(2,0)=5$，又设 $z=z(x,y)$ 是由方程 $xz=f(2x-y,xyz)$ 确定的隐函数，则 $\displaystyle \left.\frac{\partial z}{\partial x}\right|_{\substack{x=1\\ y=0}}=\underline{\hspace{3.5em}}$．
  \item 设曲线 $L$ 为球面 $x^2+y^2+z^2=4$ 与平面 $x+y+z=3$ 的交线，则曲线积分
  \[
    \int_L (x^2+2xy)\dd s=\underline{\hspace{10em}}.
  \]
\end{enumerate}
\begin{solution}
  \textbf{(1)} 根据 Taylor 公式，
  \[
    \ln(1+x)=x-\frac{x^2}{2}+\frac{x^3}{3}+o(x^3),\qquad
    \sin x=x-\frac{x^3}{6}+o(x^3),
  \]
  所以
  \[
    f(x)=-\frac{1}{3}x^3+o(x^3).
  \]
  因为 $\displaystyle \lim_{x\to 0}\frac{f(x)}{g(x)}=1$，即
  \[
    \lim_{x\to 0}\frac{-\frac{1}{3}x^3+o(x^3)}{ax^3}=1,
  \]
  所以 $a=-\frac{1}{3}$．

  \textbf{(2)} 易得此平面方程为 $2x+2y-3z=0$．

  注：答案不唯一，系数差个倍数都对，例如：$-2x-2y+3z=0$ 或 $x+y-\frac{3}{2}z=0$ 也对．

  \textbf{(3)}
  \begin{align*}
    \lim_{x\to 0}\frac{\displaystyle\int_0^{\sin^2 x}\ln(1+t)\dd t}{(\sqrt[3]{1+x^3}-1)\sin x}
    &=\lim_{x\to 0}\frac{\displaystyle\int_0^{\sin^2 x}\ln(1+t)\dd t}{\frac{1}{3}x^4}\\
    &=\lim_{x\to 0}\frac{\ln(1+\sin^2 x)\cdot 2\sin x\cos x}{\frac{4}{3}x^3}
    =\frac{3}{2}.
  \end{align*}

  \textbf{(4)} 对方程两边关于 $x$ 求偏导数，得
  \[
    z+x\frac{\partial z}{\partial x}=2f_u'+yf_v'\left(z+x\frac{\partial z}{\partial x}\right).
  \]
  由此解得
  \[
    \frac{\partial z}{\partial x}=\frac{1}{x}\left(\frac{2f_u'}{1-yf_v'}-z\right).
  \]
  注意到 $x=1$，$y=0$ 时，$z=f(2,0)=3$，$f_u'(2,0)=5$，所以
  \[
    \left.\frac{\partial z}{\partial x}\right|_{\substack{x=1\\ y=0}}=2f_u'(2,0)-3=7.
  \]

  \textbf{(5)} 根据对称性，得
  \begin{align*}
    \int_L (x^2+2xy)\dd s
    &=\frac{1}{3}\int_L (x^2+y^2+z^2)\dd s+\frac{2}{3}\int_L (xy+yz+zx)\dd s\\
    &=\frac{1}{3}\int_L (x^2+y^2+z^2)\dd s
      +\frac{1}{3}\int_L \bigl((x+y+z)^2-(x^2+y^2+z^2)\bigr)\dd s\\
    &=\frac{1}{3}\int_L (x+y+z)^2\dd s=6\pi.
  \end{align*}
\end{solution}

\answerproblem{二、}{14}
设 $f(x)$ 在 $(-\infty,+\infty)$ 上连续，在 $x=0$ 处可导，且满足
\[
  f(x)=x^3+x^2\lim_{x\to 0}\frac{f(x)}{x}-x\int_0^1 f(x)\dd x.
\]
求函数 $f(x)$ 的表达式．
\begin{solution}
  易知 $f(0)=0$，所以 $\displaystyle \lim_{x\to 0}\frac{f(x)}{x}=\lim_{x\to 0}\frac{f(x)-f(0)}{x-0}=f'(0)$．

  记 $a=f'(0)$，$b=\displaystyle\int_0^1 f(x)\dd x$，则
  \begin{equation}
    f(x)=x^3+ax^2-bx. \tag{1}
  \end{equation}
  由此得 $a=\displaystyle \lim_{x\to 0}\frac{f(x)}{x}=-b$．

  再对式子 (1) 的两边在 $[0,1]$ 上作定积分，得
  \begin{align*}
    b&=\int_0^1 x^3\dd x+a\int_0^1 x^2\dd x-b\int_0^1 x\dd x\\
    &=\frac{1}{4}+\frac{1}{3}a-\frac{1}{2}b.
  \end{align*}
  与 $a=-b$ 联立解得 $a=-\frac{3}{22}$，$b=\frac{3}{22}$，所以
  \[
    f(x)=x^3-\frac{3}{22}x^2-\frac{3}{22}x,\qquad -\infty<x<+\infty.
  \]
\end{solution}

\answerproblem{三、}{14}
求微分方程 $y'x\ln x\sin y+(1-x\cos y)\cos y=0$ 的通解．
\begin{solution}
  将原方程变形为
  \begin{equation}
    y'\sin y+\frac{1}{x\ln x}\cos y=\frac{1}{\ln x}\cos^2 y. \tag{1}
  \end{equation}
  令 $z=\cos y$，则上述方程可化为 $n=2$ 的 Bernoulli 方程
  \[
    z'-\frac{1}{x\ln x}z=-\frac{1}{\ln x}z^2.
  \]
  令 $u=z^{-1}$，原方程可进一步化为非齐次线性方程
  \[
    u'+\frac{1}{x\ln x}u=\frac{1}{\ln x}.
  \]
  利用公式，得上述方程的通解为
  \[
    u=\ee^{-\int\frac{\dd x}{x\ln x}}
    \left(\int\frac{1}{\ln x}\ee^{\int\frac{\dd x}{x\ln x}}\dd x+C\right)
    =\frac{1}{\ln x}(x+C).
  \]
  将其中的变量换回原变量：$u=z^{-1}=\frac{1}{\cos y}$，得所求通解为
  \[
    \frac{1}{\cos y}=\frac{1}{\ln x}(x+C).
  \]
\end{solution}

\answerproblem{四、}{14}
设 $\Omega$ 是由平面 $\Pi_1:\ x=0$，$\Pi_2:\ x=y$，$\Pi_3:\ y=z$，$\Pi_4:\ z=1$，$\Pi_5:\ z=2$ 围成的空间区域，求三重积分
\[
  \iiint_{\Omega}\frac{\dd x\dd y\dd z}{\sqrt{y^2+z^2}}.
\]
\begin{solution}
  采用先二后一法．设 $1\le z\le 2$，用平面 $z=z$ 截 $\Omega$ 得平面区域：$D_z=\{(x,y)\mid 0\le x\le z,\ 0\le y\le x\}$，则
  \begin{align*}
    \iiint_{\Omega}\frac{\dd x\dd y\dd z}{\sqrt{y^2+z^2}}
    &=\int_1^2 \dd z\iint_{D_z}\frac{\dd x\dd y}{\sqrt{y^2+z^2}}\\
    &=\int_1^2 \dd z\int_0^z \dd y\int_0^y \frac{\dd x}{\sqrt{y^2+z^2}}\\
    &=\int_1^2 \dd z\int_0^z \frac{y}{\sqrt{y^2+z^2}}\dd y\\
    &=\int_1^2 \left.\sqrt{y^2+z^2}\right|_0^z \dd z
      =(\sqrt{2}-1)\int_1^2 z\dd z\\
    &=\frac{3(\sqrt{2}-1)}{2}.
  \end{align*}
\end{solution}

\answerproblem{五、}{14}
判断级数 $\displaystyle \sum_{n=1}^{\infty}a_n$ 的收敛性，其中 $a_n=\sin^2 x\sin^2 2x\cdots\sin^2 2^n x$，$x\in(-\infty,+\infty)$．
\begin{solution}
  \textbf{(1)} 令 $f(x)=\sin^2 x\sin 2x$，则
  \[
    f'(x)=2\sin^2 x(3-4\sin^2 x)
    =-8\sin^2 x\left(\sin x-\frac{\sqrt{3}}{2}\right)\left(\sin x+\frac{\sqrt{3}}{2}\right).
  \]
  令 $f'(x)=0$，在 $(0,\pi)$ 中得到两个根 $x_1=\frac{\pi}{3}$，$x_2=\frac{2\pi}{3}$．

  容易证明函数 $f(x)$ 在 $\left(0,\frac{\pi}{3}\right)\cup\left(\frac{2\pi}{3},\pi\right)$ 上单调增加，在 $\left(\frac{\pi}{3},\frac{2\pi}{3}\right)$ 上单调减少．

  \textbf{(2)} 由于 $f(x)$ 的周期为 $\pi$，$f(0)=f(\pi)=0$，$f(x)$ 在 $[0,\pi]$ 上的最大值和最小值分别为 $f\left(\frac{\pi}{3}\right)=\frac{3\sqrt{3}}{8}$ 和 $f\left(\frac{2\pi}{3}\right)=-\frac{3\sqrt{3}}{8}$．故 $\abs{f(x)}\le\frac{3\sqrt{3}}{8}$．

  最后，由
  \[
    a_n^{3/2}=\abs{\sin x}\cdot\abs{f(x)f(2x)\cdots f(2^{n-1}x)}\cdot\abs{\sin^2 2^n x}
    \le\left(\frac{3\sqrt{3}}{8}\right)^n,
  \]
  知 $a_n\le\left(\frac{3\sqrt{3}}{8}\right)^{\frac{2n}{3}}=\left(\frac{3}{4}\right)^n$．从而级数 $\displaystyle \sum_{n=1}^{\infty}a_n$ 收敛．
\end{solution}

\answerproblem{六、}{14}
设函数 $f(x),g(x)$ 在 $[a,b]$ 上连续可导且满足以下三个条件：
\begin{enumerate}[label=(\roman*)]
  \item $\displaystyle \int_a^b f(x)\dd x=\int_a^b g(x)\dd x=0$；
  \item $f'(x)g(x)-f(x)g'(x)>0$，$x\in[a,b]$；
  \item $g(x)$ 在 $(a,b)$ 上有唯一零点 $\xi$．
\end{enumerate}
证明：
\begin{enumerate}
  \item $\displaystyle \int_a^b f(x)\left(\int_a^x g(s)\dd s\right)\dd x=-\int_a^b g(x)\left(\int_a^x f(s)\dd s\right)\dd x$．
  \item $\displaystyle \int_a^b f(x)\left(\int_a^x g(s)\dd s\right)\dd x>0$．
\end{enumerate}
\begin{solution}
  \textbf{(1)} 记 $F(x)=\displaystyle\int_a^x f(s)\dd s$，$G(x)=\displaystyle\int_a^x g(s)\dd s$，则由分部积分
  \[
    \int_a^b f(x)G(x)\dd x=\int_a^b G(x)\dd F(x)=-\int_a^b F(x)g(x)\dd x,
  \]
  这里，用到了条件 $F(b)=G(b)=0$．

  \textbf{(2)} 显然，$g'(\xi)\ne 0$，不妨设 $g'(\xi)>0$，则在 $[a,\xi]$ 上 $g(x)<0$，在 $(\xi,b]$ 上 $g(x)>0$．
  \[
    \int_a^b f(x)\left(\int_a^x g(s)\dd s\right)\dd x
    =\frac{1}{2}\iint_D\bigl(f(x)g(s)-f(s)g(x)\bigr)\dd x\dd s,
  \]
  其中，$D=\{(x,s)\mid a\le s\le x\le b\}$．

  把 $D$ 分成三部分：$D=D_1\cup D_2\cup D_3$，其中，$D_1=\{(x,s)\mid a\le s\le x\le \xi\}$，$D_2=\{(x,s)\mid a\le s\le\xi\le x\le b\}$，$D_3=\{(x,s)\mid \xi\le s\le x\le b\}$．

  由条件 $f'(x)g(x)-f(x)g'(x)>0$ 知 $\frac{f(x)}{g(x)}$ 在 $[a,\xi]$ 和 $[\xi,b]$ 上严格单调递增，从而在 $D_1$ 和 $D_3$ 上，$f(x)g(s)-f(s)g(x)\ge 0$ 且不恒等于 $0$，故
  \[
    \iint_{D_1\cup D_3}\bigl(f(x)g(s)-f(s)g(x)\bigr)\dd x\dd s>0.
  \]
  又
  \begin{align*}
    \iint_{D_2}\bigl(f(x)g(s)-f(s)g(x)\bigr)\dd x\dd s
    &=\int_\xi^b f(x)\dd x\int_a^\xi g(s)\dd s-\int_\xi^b g(x)\dd x\int_a^\xi f(s)\dd s\\
    &=0,
  \end{align*}
  最后一个等式用到了 $\displaystyle\int_\xi^b f(x)\dd x=-\int_a^\xi f(s)\dd s$，$\displaystyle\int_\xi^b g(x)\dd x=-\int_a^\xi g(s)\dd s$．

  综上，$\displaystyle\int_a^b f(x)\left(\int_a^x g(s)\dd s\right)\dd x>0$．
\end{solution}
